\subsection*{Additional human questions improve evidence selection}
We tested whether experience with more passages could improve the reader's choice of answer. Both fits started from the same 101,144-parameter checkpoint and saw the same number of pictured questions. One replayed the original CMRC training data, and the other replaced half of those presentations with human-authored traditional Chinese questions from DRCD. The mixed reader gave 1,270 exact stopped answers to 3,505 DRCD questions, compared with 380 under replay and 500 before further training. Against replay it gained 947 answers and lost 57, a difference of 25.39 percentage points (95\% article-bootstrap interval, 23.67 to 27.14). It also improved on official CMRC singles, from 999 before training to 1,111, and on the short instruction asking for two answers, from 505 to 601 pairs correct in both orders (Fig.~\ref{fig:traditional-passage}).

Exact image scoring combines finding a fact with selecting its boundaries. Asked in which year nearly 50,000 people in India could speak Sanskrit, the mixed reader selected 1991. The replay reader copied a longer passage that already included that date. Among the 947 newly exact DRCD answers over replay, 331 previously contained a whole accepted answer with additional ink. In 616 cases the previous output lacked a complete accepted answer or stop. The first remaining failure substituted a German scholar's name for the requested region, Europe. On the same 4,638 DRCD question pairs, the mixed reader gave both correct orders for 613 pairs with the short instruction but only 77 with the longer traditional instruction. Additional passage evidence improves factual extraction and leaves a substantial dependence on how the task is phrased.
